\section{要多少数据：Scaling law 与生命周期成本}

决定数据源之前，先要估算需要多少 tokens。Scaling law 的价值是给预算分配提供经验曲线；它的局限是训练 loss 的 optimum 不自动等于部署成本、数据受限和 domain-specific performance 的 optimum。

\subsection{三个问题：多少、哪里、怎样过滤}

上一章把模型、GPU 与数据放进同一个约束问题；本节开始拆解数据变量。读图时应把三个问句看成有依赖的 pipeline：数据量决定 collection 规模，来源决定 provenance 与 domain，filter 决定 quality/diversity；三者相互作用，不能先固定一个再把其他视为清洁步骤。

\lecturefigure{slide-13.jpg}{训练数据的三个问题：how much、where、how to filter}{Loubna Ben Allal 官方 deck，第 13 页}

\begin{importantbox}{Raw data、retained data 与 training tokens 要区分}
Raw bytes 是抓取体量；retained bytes 是 filter/dedup 后文件；unique tokens 是 tokenizer 后去重统计；training tokens 还包含 epochs/repeats 与 mixture sampling。报告“3T tokens”若不说明是哪一层，无法复现。
\end{importantbox}

\subsection{Scaling law：在参数与 tokens 之间分配 compute}

确定“需要多少数据”不能只看已有硬盘容量，而要把模型大小、训练 tokens 和预算同时放进经验曲线。本节先给出 scaling-law 形式，再用 Kaplan 与 Chinchilla 的不同结论说明：所谓 optimum 是在一组假设和 fitted range 下得到的资源配置，而不是脱离数据与硬件的常数。

Kaplan-style empirical law 常写成
\[
L(N,D)=L_\infty+aN^{-\alpha}+bD^{-\beta},
\]
其中 $N$ 为参数量、$D$ 为训练 token、$L_\infty$ 为不可约 loss。对 dense Transformer，训练 FLOPs 常近似 $C\approx 6ND$，固定 $C$ 时就需要选择模型与数据的组合。

\lecturefigure{slide-14.jpg}{How much data：用 scaling laws 分配 compute budget}{Loubna Ben Allal 官方 deck，第 14 页}

\lecturefigure{slide-15.jpg}{Kaplan 与 Chinchilla：不同 compute-optimal 结论及示例}{Loubna Ben Allal 官方 deck，第 15 页}

Kaplan 早期结果给出较大的模型、较少的 token；Chinchilla 通过更广实验修正为参数与 token 更平衡的增长。slide 对比 GPT-3 175B/约 300B tokens 与 Chinchilla 65B/约 1.4--1.6T tokens，教学结论是：\textbf{同 compute 下，旧模型可能参数过大、训练不足}。

\begin{warningbox}{常数比率不是永恒法则}
“每个参数约 20 tokens”依赖模型族、数据、optimizer、context 和 fitted range。后续模型常 intentionally overtrain，domain 或 data quality 变化也会移动 optimum。
\end{warningbox}

\teachervoice{00:06:20--00:09:10，讲者用 Kaplan/Chinchilla 说明 scaling law 是预算分配工具，并提醒这些结论来自特定论文和数据，不应只背一个 token-per-parameter 数字。}

\subsection{读 Chinchilla 曲线：每个 compute budget 都有谷底}

上一节给出了抽象公式，本节转向经验曲线，说明 optimum 在图上怎样出现。读图时先固定一种颜色：横轴是参数量，纵轴是 validation loss；每条曲线先降后升，模型太小欠容量，模型太大则在固定预算下 token 不足，谷底才是该 budget 下的 training-compute optimum。

\lecturefigure{slide-16.jpg}{Chinchilla empirical curves：固定 compute 下的参数量谷底}{Loubna Ben Allal 官方 deck，第 16 页}

若 $C=6ND$，将 $D=C/(6N)$ 代回经验 law，可以求 $N^*(C)$ 与 $D^*(C)$ 的幂律。真实训练不会只解公式，还要考虑 discrete hardware shape、parallel efficiency、batch、context 和数据可得性。

\begin{knowledgebox}{曲线不告诉你的四件事}
它通常不计 serving cost；不保证 downstream task 同步提升；不描述 license/quality constraints；也不包含训练 failure、checkpoint selection 与 hyperparameter search 成本。
\end{knowledgebox}

\subsection{Compute、Data、Model 三条 scaling evidence}

Chinchilla 曲线展示固定预算内的谷底，这一页进一步把 loss 分别对 compute、data size 和 model size 作图，展示可拟合直线区间。读图时首先确认其他变量是否接近最优，并检查横轴范围；否则单独增加某一轴可能只是把系统推入 under-training 或 data-limited regime。

\lecturefigure{slide-17.jpg}{Scaling evidence：loss 随 compute、data 与 model size 的幂律变化}{Loubna Ben Allal 官方 deck，第 17 页}

\begin{warningbox}{三张直线不能拼成无条件因果}
这些图是经验边际关系。若 data quality 随规模下降、重复率上升，或模型超出优化稳定区间，同一斜率不会无限延伸。预测必须报告 fitted range 与 confidence。
\end{warningbox}

\subsection{从 Chinchilla 到 overtraining：optimum 发生了什么变化}

时间线展示 Chinchilla 65B/1.4T、LLaMA 7B/1T、LLaMA-3 8B/15T 等趋势。团队主动让较小模型训练更多 tokens，因为一次训练的额外成本，可以换取长期更低的 inference memory、latency 与 deployment complexity。

\lecturefigure{slide-18.jpg}{Scaling-law timeline：从 compute-optimal 到小模型 overtraining}{Loubna Ben Allal 官方 deck，第 18 页}

定义模型生命周期成本
\[
C_{\text{life}}=C_{\text{train}}+Q\cdot C_{\text{infer/token}}+C_{\text{memory}}+C_{\text{ops}},
\]
其中 $Q$ 是服务期总 token。若 $Q$ 很大，多训练一个 smaller model 可能降低总成本。

\teachervoice{00:09:10--00:11:10，讲者说 Chinchilla scaling law 没把 inference cost 放入目标，所以今天会选择更小、训练更久的模型，并在训练时多付一次成本。}

\subsection{GPT-1 到 GPT-4：scale 同时放大训练与服务账单}

前面的 optimum 讨论仍然很抽象；这里用 GPT 系列的数量级变化把训练和服务账单放到同一画面。slides 用 dataset/model size 和粗略训练成本展示规模增长，数字是公开估计而不是审计账本；其教学作用是提醒读者，参数增大不仅影响一次训练，也使每 token latency、GPU memory 和并发服务长期变贵。

\lecturefigure{slide-19.jpg}{从 GPT-1 到 GPT-4：dataset、model size 与训练成本增长}{Loubna Ben Allal 官方 deck，第 19 页}

\lecturefigure{slide-20.jpg}{Inference gets more expensive：参数量、显存与 GPU 需求}{Loubna Ben Allal 官方 deck，第 20 页}

若权重精度为 $b$ bytes，单份权重 memory 至少为
\[
M_{\text{weights}}=bN,
\]
还需 KV cache、activations、workspace 与并行冗余。量化降低 $b$，但不自动降低所有 compute 和 bandwidth bottleneck。

\begin{knowledgebox}{为什么 IDE code model 偏好更小}
代码补全要求低 latency、高并发、较长 repository context，用户每次按键都可能触发请求。一个 benchmark 更强但延迟更高的模型，未必是产品 optimum；这也解释 StarCoder2 提供多种尺寸。
\end{knowledgebox}

\subsection{Compute-optimal 不等于 optimal：Harm's law 的警告}

上一节说明大模型会反复支付 inference cost；本节因此把目标从单次训练 loss 改成生命周期成本。两页 slide 写明：smaller model 训练更久会提高 upfront compute，却可能降低每次 deployment 成本；所谓 ``Harm's law'' 只是对越过 Chinchilla 点后额外训练开销的幽默命名，不是新的自然定律。

\lecturefigure{slide-21.jpg}{Compute-optimal $\neq$ deployment-optimal：为什么过训练小模型}{Loubna Ben Allal 官方 deck，第 21 页}

\lecturefigure{slide-22.jpg}{Harm's law：训练越过 Chinchilla 点的额外 compute}{Loubna Ben Allal 官方 deck，第 22 页}

\begin{importantbox}{做决策时至少画两条曲线}
一条是 training loss versus training compute，另一条是 end-to-end product cost/quality versus model size。若只优化第一条，会忽略 repeated inference、memory capacity、energy、edge deployment 与 user latency。
\end{importantbox}

\subsection{Data constraint 与 quality 会移动 scaling curve}

进一步阅读页引入 data-constrained scaling 和 DeepSeek LLM：有限 unique data 时，重复训练会改变收益；更高质量 corpus 也会改变 model/data allocation。换句话说，$D$ 不只是 token 计数，还带 distribution 和 quality。

\lecturefigure{slide-23.jpg}{Scaling laws 的进一步阅读：data constraint 与 quality-dependent allocation}{Loubna Ben Allal 官方 deck，第 23 页}

\begin{warningbox}{Domain scaling 仍是 underexplored 问题}
Q\&A 中讲者说 Chinchilla 在 English/code 间有一定支持，但 medical 等 domain 可能不同；甚至同 domain 换成更 curated data，optimum 也可能移动。不要把 generic web fit 直接复制到专门领域。
\end{warningbox}

\teachervoice{00:11:10--00:13:20，讲者把 data quality 和 domain 作为 scaling-law caveat，并指出 DeepSeek 的结果说明同一领域的数据集变化也会改变分配。}

\subsection{本章小结}

Scaling law 回答“固定训练 compute 下怎样分配”，却不回答整个模型生命周期。真正的 optimum 还取决于 serving volume、memory、latency、unique-data constraint、domain 与 quality，因此现代团队常选择更小模型加更多 tokens。

